\section{Hölder 不等式与 Minkowski 不等式}

在本小节及下一小节中，X 与 P 的含义与第 1 小节中的相同，M 表示定义在 P 上且满足第 1 小节所列条件的函数。

命题 2. --- 设 \(\alpha\) 与 \(\beta\) 是两个数，满足 \(0 < \alpha < 1\) ，\(0 < \beta < 1\) ，\(\alpha + \beta = 1\) 。若 \(f\) 与 \(g\) 是定义在 X 上的两个有限函数 \(\geqslant 0\) ，且 \(M(f)\) 与 \(M(g)\) 有限，则

\[
M (f ^ {\alpha} g ^ {\beta}) \leqslant (M (f)) ^ {\alpha} (M (g)) ^ {\beta}\tag{5}
\]

（Hölder 不等式）。

由命题 1，问题归结为证明：在 \(\mathbf{R}^2\) 中，由关系 \(t_1 \geqslant 0\) 、\(t_2 \geqslant 0\) 、\(t_1^\alpha t_2^\beta \geqslant 1\) 所定义的集是凸的，或者等价地 (FRV, I, §4, No. 1, Def. 1)，证明函数 \(u(t) = t^{-\alpha/\beta}\) 在 \(0 < t < +\infty\) 上是凸的。而若令 \(r = \alpha/\beta\) ，我们有 \(\mathrm{D}^2 u(t) = r(r + 1)t^{-r - 2}\) ，且由于 \(r > 0\) ，\(\mathrm{D}^2 u(t) > 0\) 在 \(]0, +\infty[\) 上成立，这就证明了本命题 (FRV, I, §4, No. 4, Prop. 8 的推论)。

推论. --- 设 \(\alpha_{i}\) ( \(1 \leqslant i \leqslant n\) ) 是 n 个数 \(\geqslant 0\) ，满足 \(\sum_{i=1}^{n} \alpha_{i} = 1\) ；又设 \(f_{i}\) ( \(1 \leqslant i \leqslant n\) ) 是定义在 X 上的 n 个函数 \(\geqslant 0\) ，使得 \(M(f_{i})\) 对 \(1 \leqslant i \leqslant n\) 有限。在这些条件下，

\[
M \left(f _ {1} ^ {\alpha_ {1}} f _ {2} ^ {\alpha_ {2}} \dots f _ {n} ^ {\alpha_ {n}}\right) \leqslant \left(M (f _ {1})\right) ^ {\alpha_ {1}} \left(M (f _ {2})\right) ^ {\alpha_ {2}} \dots \left(M (f _ {n})\right) ^ {\alpha_ {n}}.\tag{6}
\]

我们只需限于对每个 i 都有 \(\alpha_{i} > 0\) 的情形。只需对 n 作归纳论证，把不等式 (5) 应用于数 \(\alpha = \alpha_{1}\) 与 \(\beta = \sum_{i=2}^{n} \alpha_{i}\) ，并应用于函数 \(f = f_{1}\) 与 \(g = (f_{2}^{\alpha_{2}} f_{3}^{\alpha_{3}} \cdots f_{n}^{\alpha_{n}})^{1/\beta}\) 。

命题 3. --- 设 p 是一个实数 \(\geqslant 1\) 。若 f 与 g 是定义在 X 上的两个有限函数 \(\geqslant 0\) ，则

\[
\left(M \big ((f + g) ^ {p} \big)\right) ^ {1 / p} \leqslant \left(M (f ^ {p})\right) ^ {1 / p} + \left(M (g ^ {p})\right) ^ {1 / p}\tag{7}
\]

（Minkowski 不等式）。

我们只需限于 \(M(f^{p})\) 与 \(M(g^{p})\) 均有限的情形。由命题 1，问题归结为证明 \(R^{2}\) 中由关系 \(t_{1} \geqslant 0\) 、\(t_{2} \geqslant 0\) 、\(t_{1}^{1/p} + t_{2}^{1/p} \geqslant 1\) 所定义的集是凸的，或者等价地，证明函数 \(u(t) = (1 - t^{1/p})^{p}\) 在 \(0 \leqslant t \leqslant 1\) 上是凸的。而我们有

\[
\mathrm{D} ^ {2} u (t) = \left(1 - \frac {1}{p}\right) t ^ {1 / p - 2} (1 - t ^ {1 / p}) ^ {p - 2} \geqslant 0
\]

对 \(0 < t \leqslant 1\) 成立，由此即得本命题。

